% SPDX-FileCopyrightText: 2026 Will Cook
% SPDX-License-Identifier: CC-BY-4.0
%
% Standalone Erdős Problem Note for #68.  Context is deliberately repeated
% from the sibling notes so that this file remains independently readable.
%
% Build from the public paper directory with `make`.
\documentclass[11pt]{article}
% SPDX-FileCopyrightText: 2026 Will Cook
% SPDX-License-Identifier: CC-BY-4.0
%
% Shared preamble for the Erdős Problem Notes: the short, standalone,
% problem-owned manuscripts covering the expansion library `ErdosProblems`.
%
% One note owns one Erdős problem.  Each note repeats whatever context its own
% reader needs, so that a reader who arrives at a single problem never has to
% assemble it from the gateway paper.  Deliberate repetition across notes is the
% design, not an oversight.
%
% Source links resolve against one pinned formal-source commit, declared once
% below.  `scripts/check_problem_note_sources.py` verifies that every authored
% (file, line, declaration) triple names that exact declaration in the pinned
% snapshot, so a link cannot silently rot as later work moves lines.

\usepackage{paper-house-style}
\usepackage{array}
\usepackage{booktabs}
\usepackage{enumitem}
\setlist{topsep=0.35em,itemsep=0.2em}
\usepackage[colorlinks=true,linkcolor=linkinternal,urlcolor=linkexternal,
  citecolor=linkinternal]{hyperref}
\hypersetup{bookmarksnumbered=true,bookmarksopen=true,bookmarksopenlevel=1,
  pdfauthor={Will Cook},pdflang={en-GB}}

% ---------------------------------------------------------------------------
% Pinned formal source.  \commit names the checkpoint every link resolves
% against; the checker reads that snapshot rather than the working tree, so the
% printed coordinates stay correct however later waves move declarations.
% ---------------------------------------------------------------------------
\newcommand{\commit}{e5fd26a6d1b1a346430205eab5385b4c9565d004}
\newcommand{\commitshort}{e5fd26a6d1b1}
\newcommand{\repobase}{https://github.com/wcook04/plectis-lean-erdos249-257/blob/\commit}
\newcommand{\sourceurl}{https://github.com/wcook04/plectis-lean-erdos249-257/tree/\commit}
\newcommand{\PX}{ErdosProblems}

% Declaration names stay inside link targets.  The printed label is either a
% spaced, lower-case reading of the declaration name (\lref) or a short authored
% phrase (\lword); neither prints a module path or a raw Lean identifier.
\ExplSyntaxOn
\cs_new_protected:Npn \noteleanlabel #1
  {
    \tl_set:Nx \l_tmpa_tl { \tl_to_str:n {#1} }
    \regex_replace_all:nnN { _+ } { \c{space} } \l_tmpa_tl
    \regex_replace_all:nnN
      { ([a-z0-9])([A-Z]) }
      { \1\c{space}\2 }
      \l_tmpa_tl
    \tl_set:Nx \l_tmpa_tl { \tl_lower_case:n { \tl_use:N \l_tmpa_tl } }
    \tl_use:N \l_tmpa_tl
  }
\ExplSyntaxOff
\newcommand{\leanlabel}[1]{\noteleanlabel{#1}}
\newcommand{\lref}[3]{\href{\repobase/\PX/#1\#L#2}{\leanlabel{#3}}}
\newcommand{\lword}[4]{\href{\repobase/\PX/#1\#L#2}{#4}}
\newcommand{\lloc}[2]{\href{\repobase/\PX/#1\#L#2}{Lean source}}

% Links into a sibling library, named repository-relative.  The expansion
% library is implicit in \lref/\lword, because that is where problem-owned
% work lands.  Several problems have their headline theorem in the older
% reviewed #249/#257 corpus, and a note that could not name that library
% would have to paraphrase a checked statement instead of linking it.
\newcommand{\mref}[3]{\href{\repobase/#1\#L#2}{\leanlabel{#3}}}
\newcommand{\mword}[4]{\href{\repobase/#1\#L#2}{#4}}
\newcommand{\mloc}[2]{\href{\repobase/#1\#L#2}{Lean source}}
\emergencystretch=2em

\newtheorem{theorem}{Theorem}[section]
\newtheorem{proposition}[theorem]{Proposition}
\newtheorem{lemma}[theorem]{Lemma}
\newtheorem{corollary}[theorem]{Corollary}
\theoremstyle{definition}
\newtheorem{definition}[theorem]{Definition}
\newtheorem{example}[theorem]{Example}
\newtheorem{problem}[theorem]{Problem}
\theoremstyle{remark}
\newtheorem*{remark}{Remark}

\newcommand{\N}{\mathbb{N}}
\newcommand{\Npos}{\mathbb{N}_{>0}}
\newcommand{\Z}{\mathbb{Z}}
\newcommand{\Q}{\mathbb{Q}}
\newcommand{\R}{\mathbb{R}}
\newcommand{\lcm}{\operatorname{lcm}}
\newcommand{\ph}{\varphi}

% The series line printed under every note's title block.
\newcommand{\seriesline}{Erd\H{o}s Problem Notes}

% Production provenance.  The wording is fixed by the corpus provenance
% contract and reproduced verbatim in every manuscript in this repository, so
% the papers cannot drift into different accounts of how they were made.
\newcommand{\noteprovenance}{\provenancenote{\emph{Authorship and AI use.} The
prose, formal proofs, and software in this version were drafted and revised by
large-language-model agents under Will Cook's direction.  Cook set the
objectives and acceptance criteria, reviewed and authorised the manuscript, and
is responsible for its contents.  The agents are production tools, not authors.
See \emph{Statements and declarations}.}}

% The status paragraph every note carries, in its own words, before any result.
% Kernel checking establishes the exact Lean proposition; it does not establish
% that the proposition is the interesting one, that it is new, or that the
% Erdős problem is closed.
\newcommand{\statusboundary}{%
  \emph{Status.}  The problem treated here is open, and this note does not
  close it.  Every statement below marked as checked is a proposition that the
  pinned Lean kernel accepts from the sources this note links to, with no
  \texttt{sorry}, no added axiom, and no unchecked evaluation.  That is a claim
  about the formal statement, not about its mathematical interest, its
  novelty, or the original problem.  The unresolved obligations are named
  exactly, in their own section, and none of the finite computations,
  reductions, or no-go results here removes one of them.}

\renewcommand{\commit}{99f4bf47422abbd8757cbb22b50ba079d764d3a7}
\renewcommand{\commitshort}{99f4bf47422a}
\usepackage{longtable}
\hypersetup{
  pdftitle={Two incomparable denominator exclusions for the sum of reciprocals of n!-1 (Erdős Problem 68)},
  pdfsubject={Erdős Problem Notes: Problem 68},
  pdfkeywords={irrationality, factorial denominators, companion constant, factorial orbit, strict successor, prime-pole residue, continued fractions, Lean 4},
}

% The prose keeps compact declaration labels; the source appendix supplies the
% commit-pinned links checked by the public release gate.
\newcommand{\pdecl}[1]{\emph{\leanlabel{#1}}}
\newcommand{\Ser}{S}
\newcommand{\chn}[2]{V_{#1}(#2)}
\newcommand{\Lc}{L}

\runninghead{Erd\H{o}s \#68: two incomparable denominator exclusions}
\title{Two Incomparable Denominator Exclusions\\for $\sum_{n\ge2}(n!-1)^{-1}$}
\subtitle{Finite certificates and exact rationality criteria for Erd\H{o}s Problem~\#68}
\author{Will Cook}
\date{31 August 2026}

\begin{document}
\maketitle
\noteprovenance

\begin{problem}[Erd\H{o}s \#68]\label{res:problem}
Is
\[
  \Ser=\sum_{n\ge2}\frac1{n!-1}
\]
irrational?
\end{problem}

\begin{abstract}
Every rational representation $S=a/q$, $q>0$, of
$S=\sum_{n\ge2}(n!-1)^{-1}$ satisfies
\[
q\nmid299999!,\qquad q\ge2^{39990}>10^{12038}.
\]
The first exclusion follows from an exact strict-successor carry test;
the second from an exact continued-fraction computation.  The central
difficulty is to retain denominator information after cancellation and
then compare it with the positive tail.  We exhibit actual prime-power
cancellation and show that, at fixed factorial moment, integral channel
corrections change the residual only by an integer.  The companion
constant $S-e+2$ gives an exact factorial-orbit criterion; the remaining assertion is a
\lword{Erdos68/FactorialZeroPlateau.lean}{1090}{irrational_factorialGapSeries_iff_cofinal_strictFacTopRat_misses}{cofinal miss of the strict factorial successor}.
\end{abstract}

\section{The denominator exclusions}\label{sec:problem}

Let $S=\sum_{n\ge2}(n!-1)^{-1}$ and
\[
H_m=\sum_{n=2}^m\frac1{n!-1},\qquad
Z_m=\lfloor m!H_m\rfloor+1,\qquad
b_m=mZ_{m-1}+1-Z_m.
\]
The divisibility exclusion in the abstract rests on the implication
\begin{equation}
S=a/q,\quad q\mid(m-1)!,\quad m\ge3
\quad\Longrightarrow\quad b_m=1.\label{eq:finite-denominator-consumer}
\end{equation}
Its proof is the factorial-tail integrality argument that Han\v{c}l and
Tijdeman use for factorial series with integer coefficients
\cite[Lemma~2.1 and the following remark, p.~385]{hancl-tijdeman}; for the
denominators $n!-1$ the tail estimate is proved directly.
Indeed, $0<m!(S-H_m)<1$, by comparison with the telescoping
series $\sum_{n>m}(n-1)/n!=1/m!$.  Thus $m!S$ is the unique integer
strictly above $m!H_m$.  The same argument at $m-1$ gives
$Z_m=m!S=mZ_{m-1}$, proving \eqref{eq:finite-denominator-consumer}.
A non-unit carry at index $m$ therefore forces
\lword{Erdos68/FactorialZeroPlateau.lean}{876}{rational_denominator_ge_of_nonunit_carry}{every rational denominator to be at least $m$},
since any smaller positive $q$ would divide $(m-1)!$.
The exact value $b_{300000}\ne1$ therefore excludes every divisor of
$299999!$ as a denominator of $S$.
Equivalently, the factorial index
$\kappa(q)=\min\{r\ge1:q\mid r!\}$ must satisfy $\kappa(q)\ge300000$.
This index, rather than the largest prime factor, is the relevant measure
of factorial divisibility; compare Sondow's use of it for $e$
\cite[\S3, Theorem~1]{sondow2006}.  His approximation bound for $e$ is not
being applied to $S$.  In particular, the exclusion does not cover every
$299999$-smooth denominator: sufficiently high powers of a small prime
need not divide $299999!$.

The independent continued-fraction exclusion bounds the denominator's size.
Neither restriction implies the other: a prime between $299999$ and $599998$
satisfies the divisibility restriction and fails the size restriction,
whereas $299999!$ does the reverse.  Both conclusions are finite; they do
not decide the irrationality question posed by Erd\H{o}s
\cite[p.~102]{erdos1988} and listed in Bloom's catalogue \cite{bloom}.
Their computational certificates are specified in \S\ref{sec:finite}.

A first reading may follow the two finite certificates in
\S\ref{sec:finite}, the exact equivalence in
Theorem~\ref{res:companion-orbit-rationality-boundary}, and the obstruction
in \eqref{eq:residual-transparency}.  The channel calculations explain why
that obstruction persists; they are not premises of either finite exclusion.

\lword{Erdos68/FactorialChannelCertificate.lean}{65}{channelEvent_eq_zero_of_not_dvd}{Adjacent factorial differences affect exactly their divisor channels}.
This makes the channel equations solvable by integral triangular
elimination.  Solving them, however, does not make the residual small:
at a fixed factorial moment every correction changes it by an integer.
Sections~\ref{sec:companion-orbit}--\ref{sec:translator} identify the exact
carry test and this limitation of channel cancellation.
Section~\ref{sec:prime-pole} then shows, at the prime $139$, why a large
common denominator need not survive reduction;
Section~\ref{sec:projection} states the additional real comparison needed.

\section{The fixed companion orbit}\label{sec:companion-orbit}

The termwise identity
\[
\frac1{n!-1}=\frac1{n!}+\frac1{n!(n!-1)}
\]
replaces the changing prefixes by the factorial orbit of one fixed number,
\[
C=\sum_{n\ge2}\frac1{n!(n!-1)}.
\]
Absolute convergence gives
\begin{equation}C+(e-2)=\Ser.\label{eq:companion-decomposition}\end{equation}
The endpoint of the exponential prefix contributes residue $1$ modulo $m$;
its positive tail removes another unit on taking the floor.  This accounts
for the distinguished residue $-2$ in the following theorem.

Canonical factorial digits come from Cantor's factorial expansion
\cite{cantor1869}, in which a real number is rational exactly when its digits
vanish from some index on.  Galambos treats the rationality of Cantor series
in \cite[Ch.~II, \S2.1, pp.~21--22]{galambos1976}, and Koepf and Schmersau
prove the irrationality direction for nonterminating factorial expansions
whose digits are not eventually maximal
\cite[Example~3.2, p.~121]{koepf-schmersau}.  The following theorem is the
analogue for the companion $C=S-e+2$, whose distinguished digit is $m-2$.

\begin{theorem}[fixed companion-orbit rationality boundary]
\label{res:companion-orbit-rationality-boundary}
The following statements are equivalent:
\begin{enumerate}
\item $\Ser\in\mathbb Q$;
\item $(\lfloor m!C\rfloor+2)\bmod m=0$ for every sufficiently large $m$.
\end{enumerate}
Consequently,
\begin{equation}
 \Ser\notin\mathbb Q
 \quad\Longleftrightarrow\quad
 (\forall B)(\exists m>B)\;
   (\lfloor m!C\rfloor+2)\bmod m\ne0 .
 \label{eq:companion-cofinal}
\end{equation}
\end{theorem}

\noindent The equivalence and its cofinal form are
\href{https://github.com/wcook04/plectis-erdos/blob/f36a98bf3d3e6f65f1074e3b800e3293d5b8a51a/ErdosProblems/Erdos68/PaperCompleteExisting.lean#L29}{companion orbit boundary}.

\begin{proof}
First suppose $\Ser=a/q$ with $q>0$, and take $m$ so large that
$q\mid(m-1)!$.  Write
\[
 E_m=m!\sum_{2\le n\le m}\frac1{n!},
 \qquad
 \varepsilon_m=m!\sum_{n>m}\frac1{n!}.
\]
Then $E_m$ is an integer, $0<\varepsilon_m<1$, and $m!\Ser$ is an integer
divisible by $m$.  Multiplying \eqref{eq:companion-decomposition} by $m!$
therefore gives
\[
 \lfloor m!C\rfloor=m!\Ser-E_m-1.
\]
Every summand of $E_m$ except the endpoint $m!/m!=1$ is divisible by $m$.
Thus $E_m\equiv1\pmod m$ and
$\lfloor m!C\rfloor\equiv-2\pmod m$.

Conversely, assume the displayed congruence from some index onward.  Let
$d_m(C)=\lfloor m!C\rfloor-m\lfloor(m-1)!C\rfloor$ be the canonical
factorial digit.  Since $0\le d_m(C)<m$, for $m\ge3$ the congruence is
equivalent to
\[
 d_m(C)=m-2.
\]
Choose $N$ beyond the exceptional indices.  The canonical factorial
expansion of $C$ then has the form
\[
 C=\lfloor C\rfloor+
   \sum_{m=2}^{N}\frac{d_m(C)}{m!}+
   \sum_{m>N}\frac{m-2}{m!}.
\]
Adding $e-2=\sum_{m\ge2}1/m!$ and using the telescoping identity
\[
 \sum_{m>N}\frac{m-1}{m!}
 =\sum_{m>N}\left(\frac1{(m-1)!}-\frac1{m!}\right)
 =\frac1{N!}
\]
gives
\[
 \Ser=\lfloor C\rfloor+
   \sum_{m=2}^{N}\frac{d_m(C)+1}{m!}+\frac1{N!},
\]
which is rational.  Negating the eventual statement yields the cofinal
formulation above.
\end{proof}

\begin{remark}\label{bdry:companion-orbit-nonconcentration}
The remaining assertion is cofinal escape of this actual factorial orbit
from $-2$.  The theorem identifies the required event without supplying it.
\end{remark}

Here and below factorial digits are the floor-defined canonical digits.
The alternative representation with an eventually maximal tail is not used.
Thus ``eventually'' means at every sufficiently large integer index, whereas
``cofinally'' means at arbitrarily large indices; neither is a statement
about a finite experimental range.

\subsection{The carry and its wrap correction}\label{sec:digits}
The canonical factorial digit is
$d_m(C)=\lfloor m!C\rfloor-m\lfloor(m-1)!C\rfloor$.
For $C_m=\sum_{n=2}^m1/(n!(n!-1))$, put
\[
\delta_m=m!(C-C_m),\qquad
\sigma_m=\mathbf1_{\{\{m!C\}<\delta_m\}}.
\]
The positive tail satisfies $0<\delta_m<1/((m+1)!-1)$.
Subtracting it crosses an integer exactly when $\sigma_m=1$, so for $m\ge3$
\begin{equation}
b_m=m-1-d_m(C)+\sigma_m-m\sigma_{m-1}.
\label{eq:companion-wrap}
\end{equation}
This identity is unconditional.  Equality
$\{m!C\}=\delta_m$ gives no wrap: the shifted prefix is then an integer.
The hypothesis needed to remove the correction is
$\sigma_m=\sigma_{m-1}=0$, not merely a small positive tail.
For integer-coefficient factorial series, Han\v{c}l and Tijdeman
prove integrality of normalised rational tails
\cite[Lemma~2.1 and the following remark, p.~385]{hancl-tijdeman}, and
classify the polynomial-coefficient case
\cite[Theorem~3.1 and Corollary~3.1, pp.~390--391]{hancl-tijdeman}.
The denominators here are $n!-1$; \eqref{eq:companion-wrap} records the extra
endpoint information required after passage to the companion.

\begin{theorem}[exact carry characterisation]
\label{res:carry-characterization}\label{res:strict-successor-complete-characterization}
The following conditions are equivalent:
\[
S\notin\mathbb Q,\qquad
(\forall B)(\exists m>B)\ b_m\ne1,\qquad
(\forall B)(\exists m>B)\ m\nmid Z_m.
\]
In particular,
\lword{Erdos68/FactorialZeroPlateau.lean}{953}{irrational_factorialGapSeries_of_cofinal_nonunit_carries}{cofinal non-unit carries imply irrationality};
equivalently, the original problem is the
\lword{Erdos68/FactorialZeroPlateau.lean}{1090}{irrational_factorialGapSeries_iff_cofinal_strictFacTopRat_misses}{strict-successor criterion}.
\end{theorem}

\noindent The pair of equivalences is
\href{https://github.com/wcook04/plectis-erdos/blob/f36a98bf3d3e6f65f1074e3b800e3293d5b8a51a/ErdosProblems/Erdos68/PaperCompleteExisting.lean#L47}{carry characterisation}.
\begin{proof}
Rationality forces eventual unit carries by
\eqref{eq:finite-denominator-consumer}.  Conversely, if $b_m=1$
eventually, then $Z_m/m!$ is eventually constant.  Since
\[
H_m<Z_m/m!\le H_m+1/m!,
\]
that rational constant is $S$.  Finally writing $x=(m-1)!H_{m-1}$ gives
$b_m=m-1-\lfloor m\{x\}+1/(m!-1)\rfloor$, hence
$-1\le b_m\le m-1$.  For $m\ge3$, these bounds together with
$Z_m=mZ_{m-1}+1-b_m$, imply $m\mid Z_m$ exactly when $b_m=1$.
\end{proof}

\section{An integral basis for factorial channels}\label{sec:channels}

We seek integer combinations that cancel selected denominator
contributions without cancelling their factorial moment.  To keep track of
these two requirements, for a finite integer vector $\lambda$ define
\[
W_{d,n}=\frac{n!}{(d!)^{\lfloor n/d\rfloor}},\qquad
M(\lambda)=\sum_n\lambda_n n!,\qquad
V_d(\lambda)=\sum_n\lambda_nW_{d,n}.
\]
The
\lword{Erdos68/FactorialChannelCertificate.lean}{25}{factorial_pow_floor_dvd_factorial}{weights are integers}.
We first allow the auxiliary coordinate $e_1$,
then impose the manuscript support $n\ge2$.  This coefficient convention
introduces no term $1/(1!-1)$ into $S$.

All vector sums are finite unless explicitly called residual series.
The integer $d\ge2$ labels a channel and $n$ labels a support index;
$M$ is an integer, not an assumed positive quantity.

\begin{theorem}[divisor-channel coordinates]\label{res:divisor-channel-coordinates}
Set
\[
T_n=ne_{n-1}-e_n,\qquad
U_n=T_n-\sum_{\substack{d\mid n\\2\le d<n}}W_{d,n}U_d
\quad(n\ge2).
\]
Then
\[
M(U_n)=0,\qquad V_d(U_n)=(d!-1)\mathbf1_{d=n}.
\]
The vectors $e_1,U_2,U_3,\ldots$ form an integral basis.  Every finite vector
has the unique finite expansion
\begin{equation}
\lambda=M(\lambda)e_1+
\sum_{d\ge2}\frac{V_d(\lambda)-M(\lambda)}{d!-1}U_d.
\label{eq:channel-basis-expansion}
\end{equation}
\end{theorem}

\noindent The observables, the unique expansion and the coefficient formula are
\href{https://github.com/wcook04/plectis-erdos/blob/f36a98bf3d3e6f65f1074e3b800e3293d5b8a51a/ErdosProblems/Erdos68/PaperCompleteDivisorCoordinates.lean#L260}{divisor channel coordinates}.
\begin{proof}
The moment of $T_n$ is zero.  The floor in $W_{d,n}$ changes between
$n-1$ and $n$ precisely when $d\mid n$, so
\[
V_d(T_n)=(d!-1)W_{d,n}\mathbf1_{d\mid n}.
\]
The recursion cancels the proper divisor channels, proving the two identities
by induction.  The vectors $e_1,T_2,\ldots,T_N$ form an integral triangular
basis with final coefficients $1,-1,\ldots,-1$.  The change from $T_n$ to
$U_n$ is integral triangular with unit diagonal.  Hence $e_1,U_2,\ldots,U_N$
is also an integral basis.  Applying $M$ and each $V_d$ identifies its
coefficients as in \eqref{eq:channel-basis-expansion}.  For $d$ beyond the
support, $V_d=M$, so no infinite sum is required.

In particular, integrality of the displayed coefficients is a consequence
of the basis calculation, not an extra divisibility assumption.
\end{proof}

The prime translator is the case $U_p=T_p$ for a prime $p\ge3$.
The recursion also isolates composite channels; for example
\begin{equation}
U_9=9e_8-e_9-5040e_2+1680e_3.
\label{res:translator}
\end{equation}
It has moment zero and only channel $9$ is nonzero.  Thus each channel can
be adjusted separately by a multiple of its own denominator.  The expansion
is \href{https://github.com/wcook04/plectis-erdos/blob/f36a98bf3d3e6f65f1074e3b800e3293d5b8a51a/ErdosProblems/Erdos68/PaperCompleteSupplementary.lean#L17}{composite translator nine} and the observable
profile is \href{https://github.com/wcook04/plectis-erdos/blob/f36a98bf3d3e6f65f1074e3b800e3293d5b8a51a/ErdosProblems/Erdos68/PaperCompleteSupplementary.lean#L45}{composite translator nine profile}.

\subsection{The support restriction and the residual}

Fix an integer $D\ge2$.  We first solve the channel equations in the
auxiliary lattice and then impose the missing $e_1$ coordinate.  These are
two separate operations.
Write
\[
L_D=\operatorname{lcm}_{2\le d\le D}(d!-1),\qquad
K_D=L_De_1-\sum_{d=2}^D\frac{L_D}{d!-1}U_d.
\]
Equation~\eqref{eq:channel-basis-expansion} implies both the channel congruence
\begin{equation}
V_d(\lambda)\equiv M(\lambda)\pmod{d!-1}
\label{res:congruence}
\end{equation}
in the equivalent integer-divisibility form
\href{https://github.com/wcook04/plectis-erdos/blob/f36a98bf3d3e6f65f1074e3b800e3293d5b8a51a/ErdosProblems/Erdos68/PaperCompleteSupplementary.lean#L53}{supported channel congruence},
and the unique description of every low-channel kernel:
\begin{equation}
V_2=\cdots=V_D=0
\quad\Longleftrightarrow\quad
\lambda=tK_D+\sum_{n>D}z_nU_n,\qquad M=tL_D.
\label{eq:low-channel-classification}
\end{equation}
For the full residual
\[
\mathcal R(\lambda)=\sum_{d\ge2}\frac{V_d(\lambda)}{d!-1},
\]
we have $\mathcal R(e_1)=S$ and $\mathcal R(U_n)=1$.  Therefore
\begin{equation}
\mathcal R\left(tK_D+\sum_{n>D}z_nU_n\right)
=tL_D(S-H_D)+\sum_{n>D}z_n.
\label{eq:residual-transparency}
\end{equation}
The series converges because $V_d=M$ beyond the support.  This equation
identifies the limit of channel adjustment: every zero-moment correction
changes the residual by an integer.  A further fractional cancellation
coordinate does not arise from these corrections.

Only finitely many $z_n$ are nonzero.  Put $a_D=[e_1]K_D$ and
$u_n=[e_1]U_n$.  Admissibility on $n\ge2$ becomes the single integer equation
\begin{equation}
t a_D+\sum_{n>D}z_nu_n=0.\label{eq:support-equation}
\end{equation}
Consequently the attainable moments are
\begin{equation}
L_D\frac{g_D}{\gcd(g_D,a_D)}\mathbb Z,
\qquad g_D=\gcd\{u_n:n>D\}.
\label{eq:attainable-moment-ideal}
\end{equation}
Indeed, finite integer combinations of the $u_n$ form $g_D\mathbb Z$.
The positive generator is attained because a gcd of any integer family is
a finite integer combination of its members.  An attaining vector is
primitive: division by a common coefficient factor greater than one would
preserve all zero channels and produce a smaller positive moment.  Odd $u_n$ vanish by induction.  For a prime $p$, the sole nonzero
proper-divisor contribution at $2p$ gives $u_{2p}=-2(2p)!/2^p\ne0$.
Every tail thus contains a nonzero coefficient, so $g_D>0$.

\paragraph{A finite certificate for the infinite moment ideal.}
The scalar recurrence is
\[
u_2=2,\qquad
u_n=-\sum_{\substack{d\mid n\\2\le d<n}}W_{d,n}u_d\quad(n>2).
\]
Once a divisor $g$ controls the computed tail, only the terms with
$d\le D$ can supply a nonzero residue to a later coefficient.
Those finitely many sources have a factorial cutoff:
\[
k!\mid W_{d,dk}=\frac{(dk)!}{(d!)^k}.
\]
The quotient by $k!$ counts partitions into $k$ unordered blocks of size $d$.

\begin{theorem}[finite determination of the moment ideal]
\label{res:finite-channel-moment-certificate}
Choose a prime $\ell$ with $D/2<\ell\le D$ and put $H=D(2\ell-1)$.
Then
\[
g_D=\gcd(u_{D+1},\ldots,u_H),\qquad H<2D^2.
\]
\end{theorem}

\noindent The statement is
\href{https://github.com/wcook04/plectis-erdos/blob/f36a98bf3d3e6f65f1074e3b800e3293d5b8a51a/ErdosProblems/Erdos68/PaperCompleteMomentHorizon.lean#L252}{finite channel moment certificate}.
\begin{proof}
Since $D<2\ell\le D(2\ell-1)=H$, the finite interval includes the nonzero
term $u_{2\ell}=-(2\ell)!/2^{\ell-1}$.  Thus its gcd $g$ is positive and
$g\mid(2\ell)!$.
For $n>H$ and a divisor $d\le D$, the integer $n/d$ is at least $2\ell$.
Hence $g\mid(n/d)!\mid W_{d,n}$.  Every early source is divisible by $g$.
Strong induction handles the remaining divisors $D<d<n$, whose coefficients
are already divisible by $g$.  Thus $g$ divides the entire tail.
Bertrand's postulate supplies $\ell$ for $D\ge3$; take $\ell=2$ for $D=2$.
Finally $H\le D(2D-1)<2D^2$.

This is a quadratic upper horizon, not an assertion that the initial block
through $2D$ always determines the tail gcd.
\end{proof}

\begin{theorem}[quotient-band breakpoint]
\label{res:bandbreakpoint}
Let $\lambda$ be a finitely supported integer vector, let $d\ge2$ and
$k\ge0$ be integers, and suppose each index $n$ in its support satisfies
$kd\le n<(k+1)d$.  Then
\[
  M(\lambda)=(d!)^k V_d(\lambda).
\]
In particular, support in $[d,2d)$ and $V_d(\lambda)=0$ force
$M(\lambda)=0$.  If all supported indices are at least $d$,
$V_d(\lambda)=0$ and $M(\lambda)\ne0$, some supported index is at least $2d$.
\end{theorem}
\begin{proof}
On the stated band $\lfloor n/d\rfloor=k$, so
$n!=(d!)^kW_{d,n}$.  Sum against $\lambda_n$; the final assertion follows
by contraposition from the first-band case.
\end{proof}

\noindent The exact factorisation is checked for every quotient band at
\lword{Erdos68/ChannelBreakpointRigidity.lean}{71}{factorialMoment_eq_factorial_pow_mul_channelNumerator_band}{the band identity},
and its zero-channel consequence is
\lword{Erdos68/ChannelBreakpointRigidity.lean}{91}{factorialMoment_eq_zero_of_channelNumerator_eq_zero_band}{band cancellation}.
The first-band form is explicit at
\lword{Erdos68/ChannelBreakpointRigidity.lean}{101}{factorialMoment_eq_factorial_mul_channelNumerator_firstBand}{first-band factorisation};
the final breakpoint alternative is
\lword{Erdos68/ChannelBreakpointRigidity.lean}{130}{exists_index_ge_two_mul_of_factorialMoment_ne_zero_of_channel_eq_zero}{breakpoint witness}.

The earlier envelope $\Lambda_n=\operatorname{lcm}(1,\ldots,n)\mid u_n$
\label{eq:channel-lcm-envelope}
remains useful for other certificates.  Finite-source checks can end sooner
because only the sum of the sources has to vanish modulo $g$.
The even-coefficient sign pattern supplies a nonzero term in every tail.
Together with \eqref{eq:attainable-moment-ideal} this determines the
unrestricted minimum positive moment for every fixed depth.  For example,
$L_4=115$, $a_4=-55$, $u_6=-180$, $u_8=-4200$ and
$23u_6-u_8=60\mid\Lambda_9$ give minimum moment $1380$, attained by
$12K_4+253U_6-11U_8$.  On support through $6$ the minimum is instead $4140$.
Minimum moment and minimum upper support are different objectives.

The first minimum allows arbitrary finite upper support; the second fixes
an additional support cap.  Neither optimisation can be substituted for the
other in a residual estimate.
Neither computation supplies a cofinal real-residual gap, and neither
alters the two finite denominator exclusions.

The problem is not reduced to a finite search by the finite horizon:
that horizon certifies the algebraic moment ideal at one fixed depth,
whereas irrationality still requires an unbounded real-residual argument.

\section{The moment cost and the choice of support}\label{sec:translator}

On the actual support $n\ge2$, one has the stronger congruence
$12\mid M-2V_2$.  For $n=2,3$, $n!=2W_{2,n}$.  For $n\ge4$, both
$n!$ and $2W_{2,n}$ are divisible by 12: writing $n=2k$ or $2k+1$,
$W_{2,2k}=k!\prod_{j=1}^k(2j-1)$ is divisible by 6 for $k\ge2$.
Since every $d!-1$ is coprime to 6, low-channel annihilation implies
\begin{equation}
12L_D\mid M.\label{eq:twelve-lcm}
\end{equation}
The factor 12 is sharp at $D=2$ and $D=3$, witnessed respectively by
$-6e_2+e_4$ and $-6e_2-8e_3+5e_4$.

The factor $L_D$ and the factor $12$ combine because
$\gcd(L_D,12)=1$; separate divisibility by two integers would not in general
imply divisibility by their product.

\subsection{Growth of the compulsory factor}
The following elementary estimate quantifies this cost:
\begin{equation}
\liminf_{N\to\infty}\frac{\log L_N}{N^{3/2}\log N}
\ge\frac{2\sqrt2}{3}.\label{res:lcm-growth}
\end{equation}
The LCM asymptotic is proved here by an ordinary mathematical argument, and
Lean checks it in \texttt{wcook04/plectis-erdos} as
\href{https://github.com/wcook04/plectis-erdos/blob/0b500c7cf8e8bb7ae343484378df02f277fb8194/lean/ErdosProblems/Erdos68/PaperCompleteLiminf.lean#L42}{common denominator growth liminf}.
The separate public Lean release
\texttt{wcook04/plectis-erdos-lean}
at commit \texttt{52f29ad1} states the bound~\eqref{res:lcm-growth}, with the
least common multiple of $n!-1$ over $2\le n\le N$ as the common denominator,
as \texttt{common\_denominator\_growth} and in extended real form as
\texttt{common\_denominator\_growth\_liminf}.  The statements are in the
Comparator challenge file
\href{https://github.com/wcook04/plectis-erdos-lean/blob/52f29ad173b04e3bac941b3663f2b9aebe5de0bb/PalomarCorpus/E68/Challenge.lean#L249}{\nolinkurl{PalomarCorpus/E68/Challenge.lean}},
where the proofs are left as \texttt{sorry}.  The solution file
\href{https://github.com/wcook04/plectis-erdos-lean/blob/52f29ad173b04e3bac941b3663f2b9aebe5de0bb/Solutions/PalomarCorpus/E68/CommonDenominatorGrowth.lean#L17}{\nolinkurl{Solutions/PalomarCorpus/E68/CommonDenominatorGrowth.lean}}
closes both by transfer from \texttt{common\_denominator\_growth} in
\nolinkurl{ErdosProblems/Erdos68/PaperCompleteAsymptotics.lean} and
\texttt{common\_denominator\_growth\_liminf} in
\nolinkurl{ErdosProblems/Erdos68/PaperCompleteLiminf.lean} of the same release.
That repository's
Linux replay of Palomar's Comparator stage accepted the entry with the Lean
kernel and with NanoDa
(\href{https://github.com/wcook04/plectis-erdos-lean/actions/runs/34782407633}{run 34782407633}).
The checked bound concerns a common denominator of the partial sums and bounds
no reduced denominator of a rational representation of the series.
The finite-block inequality has a separate
Lean source,
\href{https://github.com/wcook04/plectis-erdos/blob/d788dd4b8c59f2246000f2ed98fffb8a5e8ac72e/lean/ErdosProblems/Erdos68/ChannelIntegralCongruence.lean#L524-L529}{finite terminal-block inequality}.
The divisibility $\gcd(i!-1,j!-1)\mid j!/i!-1$ follows on multiplying $i!-1$
by $j!/i!$ and subtracting $j!-1$.  Shared prime-power divisibility was
already used to control the spacing of factorial indices by Erd\H{o}s and
Stewart \cite[\S3, pp.~516--517]{erdos-stewart1976}.  The displayed
subtraction is the case $P=-1$ of Luca and Shparlinski's polynomial-shift
argument \cite[proof of Lemma~5, p.~811]{luca-shparlinski}, also used by Lai
\cite[proof of Lemma~2.4, display~(2.5)]{lai}.  These sources supply the
inherited collision step, not the lcm asymptotic stated here.
For positive integers $x_i$, the divisibility
$\prod_i x_i\mid\operatorname{lcm}(x_i)\prod_{i<j}\gcd(x_i,x_j)$
follows prime by prime: after ordering valuations, each except the maximum
occurs at least once among the pairwise minima.  On a terminal block of $k$
terms, this inequality and
$\gcd(i!-1,j!-1)\mid j!/i!-1$ give
\[
\log L_N\ge\sum_{n=N-k+1}^N\log(n!-1)
-\binom{k+1}{3}\log N.
\]
Taking $k=\lfloor\alpha\sqrt N\rfloor$ gives leading coefficient
$\alpha-\alpha^3/6$, maximised by $\alpha=\sqrt2$.  For a fixed
$P\in\mathbb Z[X]\setminus\{0\}$, the same ordinary argument gives this
bound for $\operatorname{lcm}_{n_0\le n\le N}(n!+P(n))$.  Here $n_0$ is
chosen so that every factor exceeds $1$ and every collision difference
below is nonzero.  Such a cutoff exists by~\cite[Lemma~2.1]{lai}, where Lai
adds the bound $n!+P(n)>1$ to the nonvanishing statement of Lemma~3 of
Luca and Shparlinski~\cite{luca-shparlinski}.  The lower cutoff matters:
$P=-2$ makes the factor at $n=2$ vanish.  Multiplying the factor at $i$ by
$j!/i!$ and subtracting the factor at $j$ gives
\[
\gcd(i!+P(i),j!+P(j))\mid P(j)-(j!/i!)P(i).
\]
For $n_0\le i<j\le N$ the
right side is a nonzero integer of absolute value $O_P(N^{\,j-i+\deg P})$,
so the block loses an additional $O_P(k^2\log N)$, which is lower order.  For
$P=0$, the lcm is only $N!$.
The non-polynomial perturbation $n!+2^n-1$ requires a different elimination
argument involving three hits and multiplicative order
\cite[Lemmas~2.1--2.3]{luca-shparlinski-exp}; it is not covered by the
polynomial-shift proof.
These are common-denominator estimates.  For $N\ge3$, any common divisor
of $(N-1)!-1$ and $N!-1$ divides $N-1$ by subtraction, and is coprime to
$N-1$ because $N-1\mid(N-1)!$.  Thus the two denominators are coprime.
Since $S-H_N>1/((N+1)!-1)$,
\[
 L_N(S-H_N)>
 \frac{((N-1)!-1)(N!-1)}{(N+1)!-1}\longrightarrow\infty.
\]
Even two terms therefore obstruct clearing by the full common denominator;
\eqref{res:lcm-growth} gives the stronger quantitative cost.
Spacing methods for factorial congruences
\cite[Lemma~2]{stewart2004} and the fixed-value multiplicity bound
\cite[arXiv v1, Theorem~12, p.~16]{garaev-luca-shparlinski}
are relevant comparisons, not ingredients of this proof.

\subsection{Remote primitive kernels}\label{sec:compressed-kernel}
Support location imposes a different requirement from low moment.
Let $D,r\ge2$ and let $L$ be a positive multiple of $2,\ldots,D$.  Put
\[
i_j=r+jL\ (0\le j<D),\quad N=r+(D-1)L,\quad
\alpha_d=(d!)^{L/d},
\]
\[
H(X)=\prod_{d=2}^D(\alpha_dX-1)=\sum_{j=0}^{D-1}h_jX^j,
\qquad A=\prod_{d=2}^D\alpha_d.
\]
The primitive kernel on this grid is
\[
\lambda_j^*=\frac{N!h_j}{A i_j!},\qquad
M=N!\prod_{d=2}^D(1-1/\alpha_d)>0.
\]
Its final coefficient is 1.

Consequently its integer coefficients have gcd one.  This explicit
primitivity does not establish that it has the smallest possible support
among all low-channel kernels.  The channel equations are polynomial vanishing
at the distinct points $\alpha_d^{-1}$.  Integrality follows by expanding
$h_j/A$ into products of reciprocal $\alpha_d$: each product divides the
factorial quotient $N!/i_j!$.  The minimum step is
$L=\operatorname{lcm}(2,\ldots,D)$.

The stronger divisibility
\[
r!\prod_{d=2}^D(L/d)!\prod_{d=2}^D(\alpha_d-1)\mid M
\]
follows by counting partitions into the specified equal-sized blocks.
In particular $\lfloor\sqrt N/2\rfloor!\mid M$: if
$v=\max(r,L/2)$, then $N<4v^2$ and $v!\mid M$.
Thus these moments absorb every fixed denominator as $N\to\infty$.

This assertion means that for each fixed $q\ge1$ there is a threshold,
depending on $q$, beyond which $q\mid M$.  It is not a uniform estimate for
denominators that grow with the support.
On the factorial-grid Cramer vector the residual differs from a nonzero
determinant times $S$ by an integer: this is the
\lword{Erdos68/PrimeUnitTranslator.lean}{1559}{exists_cramerFactorialGridResidual_eq_det_mul_factorialGapSeries_add_int}{Cramer residual identity}.
Cofinal residual nonintegrality on this family is sufficient for irrationality;
the construction does not establish that nonintegrality.

\section{Cancellation in a reduced prefix}\label{sec:prime-pole}

A common-denominator bound does not determine the denominator after addition.
For linear forms in $p$-adic zeta values, with $p\ge5$ a fixed prime and $n$
the index of the form, Lai, Lupu and Sprang quantify a gap of this kind.
They bound the denominators of the coefficients by powers of
$\operatorname{lcm}(1,\ldots,n)$ and show that the bound may be divided by a
product $\Phi_n$ of prime powers~\cite[\S5]{lai-lupu-sprang}.  The growth
rate of $\Phi_n$~\cite[Lemma~7.3]{lai-lupu-sprang} enters the
inequality~\cite[(8.1)]{lai-lupu-sprang} under which the rescaled forms,
which have integer coefficients, meet the irrationality criterion
\cite[Lemma~2.1, p.~3]{lai-lupu-sprang} that they quote from Lai
\cite[Lemma~2.1, p.~4]{lai-2adic}: nonzero values whose $p$-adic size,
multiplied by the largest absolute value of a coefficient, tends to zero.
For the prefix sums no asymptotic saving of this kind is established here;
the first-layer test below decides, one prime at a time, whether the largest
power of that prime survives reduction.  This paragraph compares the two
methods; no $p$-adic theorem is applied to $S$.

The test is the top valuation layer of the reciprocal-sum valuation formula
of Louwsma and Martino \cite[Lemma~4.1, p.~10]{louwsma-martino}, written out
here for the denominators $n!-1$.
Put $L_M=\operatorname{lcm}_{2\le n\le M}(n!-1)$ and
$A_M=\sum_{n=2}^M L_M/(n!-1)$.  For a prime $q$ with attained maximum
$e=\max_{2\le n\le M}v_q(n!-1)>0$, set
\[
I_{q,e}(M)=\{n\in[2,M]:v_q(n!-1)=e\},\qquad
R_{q,e}(M)=\sum_{n\in I_{q,e}(M)}
\left(\frac{n!-1}{q^e}\right)^{-1}\pmod q.
\]
After multiplication by $L_M$, all lower-valuation terms vanish modulo $q$.
Hence
\[
A_M\equiv (L_M/q^e)R_{q,e}(M)\pmod q,
\]
and $L_M/q^e$ is a unit.  Consequently
\begin{equation}
v_q(\operatorname{den}H_M)=e
\quad\Longleftrightarrow\quad R_{q,e}(M)\ne0\pmod q.
\label{eq:prime-pole-survival}
\end{equation}
The equivalence, with the maximum and its attainment taken from the actual
prefix lcm, is \href{https://github.com/wcook04/plectis-erdos/blob/f36a98bf3d3e6f65f1074e3b800e3293d5b8a51a/ErdosProblems/Erdos68/PaperCompletePrimePole.lean#L118}{maximal prime power survival}.

The following factorial-gap examples come from an exact modular scan; they
show why the weights matter.  At
$M=138$, the maximal $139$-hits are $69,122,137$, with cofactor residues
$6,49,73$, and
\[
6^{-1}+49^{-1}+73^{-1}=0\pmod{139}.
\]
At $M=2592$, the maximal $2593$-hits are $349,2243,2591$, with residues
$1508,1566,1678$ and reciprocal sum zero.  All these attained exponents are
one, so the respective prime disappears completely from the reduced prefix.
Singleton maximal support cannot cancel, whereas a hit-count bound alone
cannot determine a weighted sum.  Even survival must subsequently be coupled
to a real strict-successor inequality.

\section{One joint collision and residue estimate}\label{sec:projection}

The whole-prefix construction separates the common collision factor from
prime-power layers with a unique maximal owner.  For a natural parameter
$p\ge3$, write $d_n=n!-1$ and put
\[
F_p=(p-1)!,\quad
D_p=\operatorname{lcm}_{2\le i<j\le2p-1}\gcd(d_i,d_j),
\]
\[
C_p=\operatorname{lcm}(F_p,D_p),\quad
L_p^{\rm blk}=\operatorname{lcm}(F_p,d_2,\ldots,d_{2p-1}),
\]
\[
R_p=L_p^{\rm blk}/C_p,\quad
T_p=\sum_{n=2}^{2p-1}L_p^{\rm blk}/d_n,\quad
\rho_p=(-T_p)\bmod R_p.
\]
Every prime dividing $R_p$ has a unique denominator attaining its maximal
valuation.  Reducing $T_p$ modulo that prime leaves one unit term, giving
$\gcd(T_p,R_p)=1$.  Thus $0<\rho_p<R_p$ whenever $R_p>1$.

\begin{proposition}[global complementary-residue criterion]
\label{res:global-complementary-criterion}
If arbitrarily large natural parameters $p\ge3$ satisfy
\begin{equation}
R_p>1,\qquad
(2p+1)L_p^{\rm blk}<2p^2(2p-1)!\rho_p,
\label{eq:global-complementary-target}
\end{equation}
then $S$ is irrational.
\end{proposition}

\noindent The criterion at natural parameters is
\href{https://github.com/wcook04/plectis-erdos/blob/f36a98bf3d3e6f65f1074e3b800e3293d5b8a51a/ErdosProblems/Erdos68/PaperCompleteExisting.lean#L156}{global complementary criterion nat}.
\begin{proof}
Suppose $S=a/q$.  For large $p$, $q\mid F_p\mid C_p$, so
$L_p^{\rm blk}S$ is an integer multiple of $R_p$.  Therefore
\[
A=L_p^{\rm blk}(S-H_{2p-1})
\]
is a positive integer congruent to $-T_p$ modulo $R_p$, and $A\ge\rho_p$.
The positive tail bound
\[
S-H_{2p-1}<\frac{2p+1}{2p^2(2p-1)!}
\]
and \eqref{eq:global-complementary-target} give $A<\rho_p$, a contradiction.
\end{proof}

Let $\widetilde C_p=C_p/F_p$ and $\eta_p=\rho_p/R_p$.  The inequality is
exactly
\begin{equation}
\log\widetilde C_p-\log\eta_p
<\log\frac{2p^2(2p-1)!}{(2p+1)(p-1)!}.
\label{eq:joint-loss-budget}
\end{equation}
The right side is $p\log p+(2\log2-1)p+O(\log p)$.
The modulus $R_p$ cancels: increasing private support alone does not improve
this comparison.  Coprimality permits $\rho_p=1$, so the complementary gap
needs its own quantitative control.  Both losses must be controlled at the
same parameters.  The
\lword{Erdos68/EndpointWeightedPrivateSupport.lean}{4766}{factorialBlock_unitFactorPairFloor_scale_iff}{unit-factor pair floor}
splits that scale test into the complementary-residue inequality and a local
collision-core cap.  A joint mean bound below $(1-\varepsilon)p\log p$ on the
same nonempty sets in $[X,2X]$ would suffice; two unrelated cofinal sets can
be disjoint.

\section{The remaining real comparison}\label{sec:open}\label{sec:plateau}

The exact target remains the
\lword{Erdos68/FactorialZeroPlateau.lean}{1090}{irrational_factorialGapSeries_iff_cofinal_strictFacTopRat_misses}{strict-successor criterion}:
\begin{equation}
(\forall B)(\exists m>B)\ m\nmid Z_m.
\label{eq:canonical-open-target}
\end{equation}
Two sufficient approaches now have explicit, different missing estimates.
For the compressed primitive grids, the moments already absorb every fixed
denominator.  With positive moment $M$ and upper support $N$, write
\[
A_N=\sum_{d=D+1}^N\frac{V_d(\lambda)}{d!-1},\qquad
\mathcal R(\lambda)=A_N+M\sum_{d>N}\frac1{d!-1}.
\]
The sufficient strict comparison is
\begin{equation}
\frac{2M}{(N+1)!-1}<\lfloor A_N\rfloor+1-A_N.
\label{eq:signed-block-gap}
\end{equation}
It puts the residual below the next integer and above $A_N$.  Proving this
on an unbounded grid family would suffice.  The finite signed block $A_N$
is essential.  The second approach asks for the simultaneous collision and
complementary-gap estimate \eqref{eq:joint-loss-budget}.
Neither estimate is established cofinally here.

In the first route $M>0$ is essential for the stated one-sided tail
comparison.  The signed prefix $A_N$ itself need not be positive.  In the
second route the core estimate and residue estimate must hold at the same
parameters; separate unbounded families do not suffice.

\paragraph{The limit of short supports.}
\label{sec:nogo}\label{sec:adjacent-unit-no-go}
Nonzero-moment kernels with $N=D+O(1)$ cannot meet
\eqref{eq:signed-block-gap}.
For $N=D+O(1)$, the compulsory divisibility $L_D\mid M$ and
\eqref{res:lcm-growth} give
\[
\frac{2|M|}{(N+1)!-1}\longrightarrow\infty,
\]
whereas the right side of \eqref{eq:signed-block-gap} is at most one.
A useful family must therefore allow a larger upper support and still
control its finite signed block.

Large denominator valuations alone also leave the real gap uncontrolled:
the model fractions $1/q^e$ tend to zero.  The exact
\lword{Erdos68/PrimeZeroBranch.lean}{6099}{factorialGap_zero_branch_iff_lower_endpoint_cylinder}{zero-branch / lower-endpoint cylinder}
and doubled-prime tests require joint numerator or predecessor information.
Their thresholds, the
\lword{Erdos68/AdjacentUnitCarryWindow.lean}{148}{consecutive_unit_carries_iff_positive_offset_le_den}{adjacent-window collapse},
fixed-owner absorption,
and the strongest countermodels are retained in the long reasoning record.

\section{The finite evidence}\label{sec:finite}

An exact non-unit carry at the prime index $67$ already forces
\lword{Erdos68/FactorialZeroPlateau.lean}{940}{sixty_seven_le_rational_denominator}{every rational denominator to be at least $67$}.
The exact GMP carry certificate covers $3\le m\le300000$.  Its unit carries
occur at
\[
52,\ 591,\ 1030,\ 1407,\ 1438,\ 2164,\ 4258,\ 10991,\ 21236.
\]
In particular $b_{300000}\ne1$, yielding $q\nmid299999!$ by
\eqref{eq:finite-denominator-consumer}.  This excludes every divisor of that
factorial, including large divisors.  It does not exclude all integers with
small prime factors, since their multiplicities can be too large.
The range through $300000$ is a fresh exact GMP replay in this revision
and lies outside the Lean development.  The long record gives the receipt's
digests and a separate replay through $4000$.  The certificate's event-trace,
unit carries and terminal enclosure match the previously retained receipt;
the published driver and backend hashes are the SPDX-bearing public copies.

The separate $80000$-bit exact rational enclosure forces $23449$ common
continued-fraction partial quotients.  The continued-fraction recurrences,
determinant identity and fractional transformation
\cite[\S1.12(ii), (1.12.5)--(1.12.7), (1.12.20)--(1.12.21)]{nist-dlmf} turn
this common prefix into the bound $q\ge2^{39990}>10^{12038}$; the enclosure,
the prefix length and the exponent are outputs of the computation.  The
open-versus-closed complete-quotient endpoint
convention matters for a sharper extraction; no stronger numerical floor is
asserted here.  The exact sources, enclosure and integer-computation receipts
are retained with the long record.  Neither finite calculation supplies
\eqref{eq:canonical-open-target}.

\appendix
\section{Guide to the formal sources}\label{app:sources}

The mathematical proofs above are independent of the source-link audit.
Links carrying the manuscript commit identify historical source locations;
they do not assert that every linked module was compiled together.
The accompanying revision bundle gives a declaration-level ledger for the
supplied public snapshot
\texttt{3d6d938d696f}: source presence, recorded build inclusion, and
computational evidence are separate fields.

The common-denominator liminf, moment horizon and moment-ideal modules have
recorded checked-build entries in that ledger.  Other supplied modules,
including \texttt{PaperCompleteExisting} and the candidate finite-size
certificate, remain outside the recorded checked build.  This revision
replayed \texttt{lake build ErdosProblems Erdos249257} on Lean~4.29.1 and
mathlib \texttt{5e932f97dd25535344f80f9dd8da3aab83df0fe6}; the library
\texttt{AxiomAudit} prints the imported paper-linked declarations against
the standard axioms \texttt{propext}, \texttt{Classical.choice}, and
\texttt{Quot.sound}.

The continued-fraction enclosure, the small arithmetic witnesses, and the
full carry census through $300000$ have fresh independent integer replays.
The earlier replay through $4000$ remains a separate smaller certificate.
The complete historical source catalogue is preserved in the handover
rather than interleaved with the mathematical argument.  Production PDF
compilation with the repository preamble, fresh formal checking and a
matching full carry receipt remain release gates.


\paragraph{Source-current statements and evidence boundaries.}

\begingroup\small\raggedright
The finite divisor-event identities, isolated-channel recursion, vanishing
moment, individual channel values, and the factor $12$ are checked in
\href{https://github.com/wcook04/plectis-erdos/blob/f36a98bf3d3e6f65f1074e3b800e3293d5b8a51a/ErdosProblems/Erdos68/DivisorChannelBasis.lean}{divisor channel basis}.  The unique integral basis expansion,
attainable-moment formula and compressed primitive-grid construction are
ordinary proofs, described in \href{https://github.com/wcook04/plectis-erdos/blob/f36a98bf3d3e6f65f1074e3b800e3293d5b8a51a/ErdosProblems/Erdos68/DivisorChannelBasis.md}{divisor channel basis} and, for the
grid construction, in
\href{https://github.com/wcook04/plectis-erdos/blob/f36a98bf3d3e6f65f1074e3b800e3293d5b8a51a/ErdosProblems/Erdos68/CompressedPrimitiveChannelKernel.md}{compressed primitive channel kernel}.  The quadratic stopping theorem
is checked in \href{https://github.com/wcook04/plectis-erdos/blob/f36a98bf3d3e6f65f1074e3b800e3293d5b8a51a/ErdosProblems/Erdos68/PaperCompleteMomentHorizon.lean}{complete moment horizon}, together with the
scalar recurrence, the equal-block divisibility and the prime anchor it uses.
The earlier lcm-envelope certificate and depth-$4$ moment $1380$ are
ordinary proofs with exact checks in \href{https://github.com/wcook04/plectis-erdos/blob/f36a98bf3d3e6f65f1074e3b800e3293d5b8a51a/ErdosProblems/Erdos68/scripts/check_moment_saturation.py}{check moment saturation}.
The finite-source checks are in \href{https://github.com/wcook04/plectis-erdos/blob/f36a98bf3d3e6f65f1074e3b800e3293d5b8a51a/ErdosProblems/Erdos68/scripts/check_finite_source.py}{check finite source}.
The inherited generic lemmas in \href{https://github.com/wcook04/plectis-erdos/blob/f36a98bf3d3e6f65f1074e3b800e3293d5b8a51a/ErdosProblems/Erdos68/TailIdealCertificate.lean}{tail ideal certificate} assume
the lcm envelope; they do not formalise the factorial arithmetic.
The growth liminf is the ordinary asymptotic consequence of the finite
terminal-block inequality cited above, and Lean checks it as
\href{https://github.com/wcook04/plectis-erdos/blob/0b500c7cf8e8bb7ae343484378df02f277fb8194/lean/ErdosProblems/Erdos68/PaperCompleteLiminf.lean#L42}{common denominator growth liminf}.
The checked counterpart of the growth liminf in the public Lean release
\texttt{wcook04/plectis-erdos-lean} at commit \texttt{52f29ad1} is cited
at~\eqref{res:lcm-growth}.

Formal proof checking does not assess external novelty or the significance
of the remaining hypothesis.  The complete source register, receipts and
failed-route calculations belong to the long record.

\par\endgroup

\paragraph{Pinned declaration index.}

The public \texttt{ErdosProblems.Erdos68} package contains the checked source
for this note.  The pinned snapshot contains fifteen cited modules:
\texttt{AdjacentUnitCarryWindow},
\texttt{CanonicalFactorialDigits},
\texttt{ChannelBreakpointRigidity},
\texttt{ChannelIntegralCongruence},
\texttt{DivisorFactorialCentre},
\texttt{EndpointWeightedPrivateSupport},
\texttt{FactorialCarry},
\texttt{FactorialChannelCertificate},
\texttt{FactorialGapPlateauCore} (isolation pin; \texttt{FactorialZeroPlateau}
is the default-root duplicate),
\texttt{FiniteDefectAutomaton},
\texttt{GapScalarNormalForm},
\texttt{PrimeThresholdParity},
\texttt{PrimeUnitTranslator},
\texttt{PrimeZeroBranch}, and
\texttt{StrictSuccessorArithmetic}.
Only these public modules belong to the manuscript source surface; no private
auxiliary digit-rigidity file is cited or projected.  The public root imports
every cited module.  The declaration table below is
pinned to the shared formal-source commit used throughout this problem-note
series.

\begin{itemize}[leftmargin=*]
\item \lword{Erdos68/AdjacentUnitCarryWindow.lean}{148}{consecutive_unit_carries_iff_positive_offset_le_den}{adjacent-window collapse}
\item \lref{Erdos68/AdjacentUnitCarryWindow.lean}{215}{twoStep_den_mul_transitionNormalizers}
\item \lref{Erdos68/AdjacentUnitCarryWindow.lean}{243}{adjacentUnitCarryWindowDen_eq_twoStep_den}
\item \lref{Erdos68/AdjacentUnitCarryWindow.lean}{337}{adjacentUnitCarryWindowOffset_eq_twoStep_reducedNumerator}
\item \lword{Erdos68/FactorialZeroPlateau.lean}{876}{rational_denominator_ge_of_nonunit_carry}{denominator at least the carry index}
\item \lword{Erdos68/FactorialZeroPlateau.lean}{940}{sixty_seven_le_rational_denominator}{denominator at least $67$}
\item \lword{Erdos68/FactorialZeroPlateau.lean}{1090}{irrational_factorialGapSeries_iff_cofinal_strictFacTopRat_misses}{strict-successor criterion}
\item \lword{Erdos68/FactorialZeroPlateau.lean}{953}{irrational_factorialGapSeries_of_cofinal_nonunit_carries}{cofinal non-unit carries imply irrationality}
\item \lword{Erdos68/PrimeZeroBranch.lean}{6099}{factorialGap_zero_branch_iff_lower_endpoint_cylinder}{zero-branch / lower-endpoint cylinder}
\item \lword{Erdos68/EndpointWeightedPrivateSupport.lean}{4766}{factorialBlock_unitFactorPairFloor_scale_iff}{unit-factor pair floor}
\item \lword{Erdos68/PrimeUnitTranslator.lean}{1559}{exists_cramerFactorialGridResidual_eq_det_mul_factorialGapSeries_add_int}{Cramer residual identity}
\item \lref{Erdos68/DivisorFactorialCentre.lean}{34}{residualCentreTerm}
\item \lref{Erdos68/DivisorFactorialCentre.lean}{39}{factorialCoeffTerm}
\item \lref{Erdos68/DivisorFactorialCentre.lean}{44}{residualCentre}
\item \lref{Erdos68/DivisorFactorialCentre.lean}{49}{factorialCoeffRat}
\item \lref{Erdos68/DivisorFactorialCentre.lean}{80}{factorial_eq_mul_pred_factorial}
\item \lref{Erdos68/DivisorFactorialCentre.lean}{88}{pred_div_eq_of_not_dvd}
\item \lref{Erdos68/DivisorFactorialCentre.lean}{102}{pred_div_add_one_eq_of_dvd}
\item \lref{Erdos68/DivisorFactorialCentre.lean}{129}{residualCentreTerm_of_not_dvd}
\item \lref{Erdos68/DivisorFactorialCentre.lean}{146}{residualCentreTerm_of_dvd}
\item \lref{Erdos68/DivisorFactorialCentre.lean}{170}{residualCentreTerm_self}
\item \lref{Erdos68/DivisorFactorialCentre.lean}{179}{factorialCoeffTerm_self}
\item \lref{Erdos68/DivisorFactorialCentre.lean}{188}{residualCentre_recurrence}
\item \lword{Erdos68/FactorialChannelCertificate.lean}{25}{factorial_pow_floor_dvd_factorial}{weights are integers}
\item \lref{Erdos68/FactorialChannelCertificate.lean}{39}{channelWeight}
\item \lref{Erdos68/FactorialChannelCertificate.lean}{43}{channelWeight_mul_denominator}
\item \lref{Erdos68/FactorialChannelCertificate.lean}{51}{channelNumerator}
\item \lref{Erdos68/FactorialChannelCertificate.lean}{55}{factorialMoment}
\item \lref{Erdos68/FactorialChannelCertificate.lean}{59}{channelEvent}
\item \lword{Erdos68/FactorialChannelCertificate.lean}{65}{channelEvent_eq_zero_of_not_dvd}{adjacent differences affect only divisor channels}
\item \href{https://github.com/wcook04/plectis-erdos/blob/99f4bf47422abbd8757cbb22b50ba079d764d3a7/ErdosProblems/Erdos68/FactorialChannelCertificate.lean#L101}{\leanlabel{lambda34}}
\item \lref{Erdos68/FactorialChannelCertificate.lean}{105}{lambda34_apply_three}
\item \lref{Erdos68/FactorialChannelCertificate.lean}{109}{lambda34_apply_four}
\item \lref{Erdos68/FactorialChannelCertificate.lean}{112}{lambda34_channel_two}
\item \lref{Erdos68/FactorialChannelCertificate.lean}{119}{lambda34_factorialMoment}
\item \lref{Erdos68/FactorialChannelCertificate.lean}{126}{lambda34_channel_three}
\item \lref{Erdos68/FactorialChannelCertificate.lean}{133}{lambda34_channel_four}
\item \lref{Erdos68/FactorialChannelCertificate.lean}{140}{lambda34_channel_ge_five}
\item \lref{Erdos68/FactorialChannelCertificate.lean}{151}{lambda34ResidualOfTail}
\item \lref{Erdos68/FactorialChannelCertificate.lean}{156}{lambda34_residual_bounds}
\item \lref{Erdos68/ChannelIntegralCongruence.lean}{868}{square_subsequence_radius_cubic_lower}
\item \lref{Erdos68/ChannelIntegralCongruence.lean}{885}{no_eventual_square_subsequence_cubic_upper}
\item \lref{Erdos68/ChannelIntegralCongruence.lean}{905}{not_isLittleO_square_subsequence_radius}
\item \lref{Erdos68/ChannelIntegralCongruence.lean}{1084}{square_subsequence_radius_three_halves_lower}
\item \lref{Erdos68/ChannelIntegralCongruence.lean}{1103}{no_eventual_square_subsequence_three_halves_upper}
\end{itemize}

\begingroup
\raggedright
\paragraph{Source-current companion orbit and Comparator routes.}
The Comparator isolation pin for the carry equivalences is
\href{https://github.com/wcook04/plectis-erdos/blob/f36a98bf3d3e6f65f1074e3b800e3293d5b8a51a/ErdosProblems/Erdos68/FactorialGapPlateauCore.lean}{factorial gap plateau core}.
\href{https://github.com/wcook04/plectis-erdos/blob/f36a98bf3d3e6f65f1074e3b800e3293d5b8a51a/ErdosProblems/Erdos68/FactorialZeroPlateau.lean}{Factorial zero plateau} is the default-root duplicate of that family and is not
the isolation source.
The complete infinite rationality boundary is checked in
\href{https://github.com/wcook04/plectis-erdos/blob/f36a98bf3d3e6f65f1074e3b800e3293d5b8a51a/ErdosProblems/Erdos68/CompanionOrbitRationality.lean}{companion orbit rationality}.  Its paper-facing
endpoints are
\href{https://github.com/wcook04/plectis-erdos/blob/f36a98bf3d3e6f65f1074e3b800e3293d5b8a51a/ErdosProblems/Erdos68/CompanionOrbitRationality.lean}{not irrational factorial gap series iff eventually companion floor $-2$}
and
\href{https://github.com/wcook04/plectis-erdos/blob/f36a98bf3d3e6f65f1074e3b800e3293d5b8a51a/ErdosProblems/Erdos68/CompanionOrbitRationality.lean}{irrational factorial gap series iff cofinal companion floor misses}.
The coherent Comparator composite
\href{https://github.com/wcook04/plectis-erdos/blob/f36a98bf3d3e6f65f1074e3b800e3293d5b8a51a/ErdosProblems/Erdos68/CompanionOrbitRationality.lean}{companion-orbit complete characterisation} routes to
Theorem~\ref{res:companion-orbit-rationality-boundary} and
Remark~\ref{bdry:companion-orbit-nonconcentration}.  The moving-factor
Comparator endpoints are recorded in the long reasoning record as the
moving-factor scale split, the split-factor normalised collision, and the
fixed-owner absorption boundary.  The separate public Lean release
\texttt{wcook04/plectis-erdos-lean} at commit \texttt{52f29ad1} selects those
three endpoints as
\texttt{movingPrivateFactorScaleSplit\_implies\_irrational},
\texttt{splitFactorNormalizedCollision\_implies\_irrational} and
\texttt{fixedOwnerPair\_eventually\_absorbed} in
\texttt{PalomarCorpus/E68/Challenge.lean}.  These source-current modules, Comparator
packages, and this manuscript stage require one common immutable public
checkpoint before terminal external replay or Palomar readiness is claimed;
for the eight-entry release that checkpoint is commit \texttt{52f29ad1},
whose Linux replay of Palomar's Comparator stage accepted the \texttt{E68}
entry (run 34782407633, cited at~\eqref{res:lcm-growth}).
The unconditional correction \eqref{eq:companion-wrap} is checked in
\href{https://github.com/wcook04/plectis-erdos/blob/f36a98bf3d3e6f65f1074e3b800e3293d5b8a51a/ErdosProblems/Erdos68/CompanionConstantCarryLaw.lean}{companion constant carry law}, and the pole-residue identity and the two
displayed reciprocal equalities are checked in
\href{https://github.com/wcook04/plectis-erdos/blob/f36a98bf3d3e6f65f1074e3b800e3293d5b8a51a/ErdosProblems/Erdos68/PrimePoleCriterion.lean}{prime pole criterion}.  The complete maximal-hit data come from the separate
exact modular scan; the carry census and continued-fraction certificate use
exact integer computation outside Lean.  The inherited public source pin is
\texttt{\commitshort}; source-current additions and that snapshot require one
immutable release before the links represent a replay of the whole note.
\par
\endgroup

\clearpage

\begin{thebibliography}{99}

\bibitem{erdos1988}
Paul Erd{\H o}s. \href{https://doi.org/10.1017/CBO9780511897184.009}{On the irrationality of certain series: problems and results}.
In Alan Baker (ed.), \emph{New Advances in Transcendence Theory}, Cambridge University Press (1988), pp.~102--109.

\bibitem{bloom}
Thomas F. Bloom. \href{https://www.erdosproblems.com/68}{Erd{\H o}s Problem {\#68}}.
Online resource (2026). Historical access: 28 July 2026; present-page status not reverified.

\bibitem{cantor1869}
Georg Cantor. {\"U}ber die einfachen Zahlensysteme.
\emph{Zeitschrift f{\"u}r Mathematik und Physik} \textbf{14} (1869), 121--128.

\bibitem{galambos1976}
J{\'a}nos Galambos. \href{https://doi.org/10.1007/BFb0081642}{Representations of Real Numbers by Infinite Series}.
Lecture Notes in Mathematics 502, Springer (1976).

\bibitem{koepf-schmersau}
Wolfram Koepf and Dieter Schmersau. \href{https://doi.org/10.1524/anly.2011.1094}{Irrationality of certain infinite series {II}}.
\emph{Analysis} \textbf{31} (2011), 117--124.

\bibitem{hancl-tijdeman}
Jaroslav Han{\v c}l and Robert Tijdeman. \href{https://doi.org/10.4064/aa118-4-5}{On the irrationality of factorial series}.
\emph{Acta Arithmetica} \textbf{118} (4) (2005), 383--401.

\bibitem{louwsma-martino}
Joel Louwsma and Joseph Martino. \href{https://arxiv.org/abs/2309.07280v1}{Rational numbers with odd greedy expansion of fixed length}.
Preprint (2023). arXiv:2309.07280.

\bibitem{lai}
Li Lai. \href{https://doi.org/10.1017/S0004972725100543}{On the largest prime divisor of {$n!+1$}}.
\emph{Bulletin of the Australian Mathematical Society} \textbf{113} (3) (2026), 390--403. arXiv:2103.14894.

\bibitem{luca-shparlinski}
Florian Luca and Igor E. Shparlinski. \href{https://doi.org/10.1112/S0024609305004923}{Prime divisors of shifted factorials}.
\emph{Bulletin of the London Mathematical Society} \textbf{37} (6) (2005), 809--817.

\bibitem{lai-lupu-sprang}
Li Lai, Cezar Lupu and Johannes Sprang. \href{https://doi.org/10.1007/s40687-025-00559-x}{On the irrationality of certain {$p$}-adic zeta values}.
\emph{Research in the Mathematical Sciences} \textbf{12} (4) (2025), article 77. arXiv:2505.23088.

\bibitem{lai-2adic}
Li Lai. \href{https://doi.org/10.1142/S1793042125500113}{On the irrationality of certain {$2$}-adic zeta values}.
\emph{International Journal of Number Theory} \textbf{21} (1) (2025), 207--235. arXiv:2304.00816.

\bibitem{nist-dlmf}
{NIST Digital Library of Mathematical Functions}. \href{https://dlmf.nist.gov/1.12}{Continued fractions: convergents}.
Online resource (2026).

\bibitem{erdos-stewart1976}
Paul Erd{\H o}s and Cameron L. Stewart. \href{https://doi.org/10.1112/jlms/s2-13.3.513}{On the greatest and least prime factors of {$n!+1$}}.
\emph{Journal of the London Mathematical Society} \textbf{13} (3) (1976), 513--519.

\bibitem{luca-shparlinski-exp}
Florian Luca and Igor E. Shparlinski. \href{https://doi.org/10.5802/jtnb.524}{On the largest prime factor of {$n!+2^n-1$}}.
\emph{Journal de Th{\'e}orie des Nombres de Bordeaux} \textbf{17} (3) (2005), 859--870.

\bibitem{sondow2006}
Jonathan Sondow. \href{https://arxiv.org/abs/0704.1282v2}{A geometric proof that {$e$} is irrational and a new measure of its irrationality}.
\emph{American Mathematical Monthly} \textbf{113} (7) (2006), 637--641. arXiv:0704.1282. Addendum: \emph{Amer. Math. Monthly} \textbf{114} (2007), 659; reading copy v2.

\bibitem{stewart2004}
Cameron L. Stewart. \href{https://doi.org/10.5486/PMD.2004.3190}{On the greatest and least prime factors of {$n!+1$}, {II}}.
\emph{Publicationes Mathematicae Debrecen} \textbf{65} (3--4) (2004), 461--480.

\bibitem{garaev-luca-shparlinski}
Moubariz Z. Garaev, Florian Luca and Igor E. Shparlinski. \href{https://doi.org/10.1090/S0002-9947-04-03612-8}{Character sums and congruences with {$n!$}}.
\emph{Transactions of the American Mathematical Society} \textbf{356} (12) (2004), 5089--5102. arXiv:math/0403422.

\end{thebibliography}

\companionpapercontext

\end{document}
